\documentclass{article}

\PassOptionsToPackage{numbers,sort&compress}{natbib}
\usepackage{neurips_2026}
\usepackage[utf8]{inputenc}
\usepackage[T1]{fontenc}
\usepackage{hyperref}
\usepackage{url}
\usepackage{booktabs}
\usepackage{amsfonts}
\usepackage{amsmath}
\usepackage{amssymb}
\usepackage{amsthm}
\usepackage{microtype}
\usepackage{xcolor}
\usepackage{algorithm}
\usepackage{algorithmic}

\newcommand{\R}{\mathbb{R}}
\newcommand{\softmin}{\operatorname*{softmin}}
\newtheorem{proposition}{Proposition}[section]

\title{CARAT: Class-Aware Robust Aggregation via Hidden Task Certificates for Poisoning-Resilient Federated Learning}

% \author{%
%   Anonymous Authors \\
%   Anonymous Institution \\
%   \texttt{anonymous@neurips.cc}
% }

\begin{document}

\maketitle

\begin{abstract}
Federated learning servers must aggregate client updates without observing client data, which leaves training vulnerable to both untargeted model poisoning and targeted backdoor attacks. Existing robust aggregators mainly rely on update geometry, coordinate-wise robust statistics, or alignment with a single trusted reference direction, but these signals can be brittle when benign clients are heterogeneous and adaptive attackers stay near the benign update manifold. We present \textbf{CARAT}, a class-aware robust aggregation rule based on \emph{hidden task certificates}: using server-side labeled reference data, CARAT instantiates multiple hidden, class-balanced micro-tasks and scores each client update by the loss reduction it induces across them. CARAT combines this task-aware signal with robust clipping, common-radius evaluation, coordinate-rank anomaly penalties, and capped-simplex weight optimization under a worst-task objective with temporal regularization, thereby favoring updates that remain broadly useful across hidden semantic slices while suppressing structurally extreme clients. In the strongest completed repository-backed CIFAR100/FangAttack pilot, CARAT suppresses attacker weights and substantially improves over Mean; however, because the current completed evidence is still training-only and the full clean-accuracy and ASR benchmark suite is not yet finished, these results should be read as initial empirical validation rather than a complete benchmark claim.
\end{abstract}

\section{Introduction}
\label{sec:introduction}

Federated learning (FL) enables many clients to train a shared model without centralizing their raw data \citep{mcmahan2017fedavg,bonawitz2017secureagg,kairouz2021federated}. This decentralization is attractive in privacy- and governance-sensitive settings, but it also leaves the server in a fundamentally asymmetric position: it must update the global model using client messages that it cannot directly verify. Poisoning attacks exploit exactly this asymmetry. Malicious clients can slow convergence, distort the learned decision boundary, or implant triggers that activate attacker-chosen behavior at test time \citep{bhagoji2019adversarial,fang2020local,shejwalkar2021manipulating,baruch2019alie,bagdasaryan2020backdoor,xie2020dba,wang2020tails,zhang2022neurotoxin}. In practice, the robustness of FL is therefore tightly coupled to the aggregation rule.

The central difficulty is that malicious behavior and benign heterogeneity can look superficially similar. Under non-i.i.d.\ partitions, benign clients legitimately differ in direction, norm, sparsity, and class emphasis \citep{zhao2018noniid,li2020fedprox,karimireddy2020scaffold,zhu2021noniidsurvey}. As a result, an update can appear unusual either because it is harmful or because it reflects a different local data distribution. Existing defenses often rely on signals that are too coarse to resolve this ambiguity reliably. Pairwise distances and geometric centrality can misclassify benign dispersion as adversarial deviation; coordinate-wise robust statistics suppress only marginal extremes and ignore cross-coordinate structure; and trusted-reference schemes reduce the server's auxiliary information to a single global direction, which can miss failures concentrated on particular classes or semantic regions \citep{blanchard2017krum,elmhamdi2018bulyan,yin2018median,xie2019zeno,chen2018draco,li2019rsa,pillutla2022rfa,cao2021fltrust}.

This leaves a key question: how should the server validate client updates when benign heterogeneity is intrinsic and the most harmful attacks may only reveal themselves on specific semantic slices of the task? What is needed is a validation signal that is semantically discriminative, robust to benign diversity, and difficult for attackers to satisfy by merely staying near the benign update manifold.

To address this, the paper proposes \textbf{CARAT}, \emph{Class-Aware Robust Aggregation via Hidden Task Certificates}. CARAT assumes server-side access to labeled reference data and uses that data to instantiate hidden, class-balanced micro-tasks. Each client update is then scored by the loss reduction it induces across these tasks, producing a certificate matrix that measures whether the update is broadly useful across hidden semantic slices rather than merely geometrically central. CARAT combines this task-aware signal with robust clipping, common-radius evaluation, coordinate-rank anomaly penalties, and capped-simplex weight optimization under a worst-task objective with temporal regularization. Our main contributions are summarized as follows.

\begin{itemize}
\item We propose CARAT, a class-aware robust aggregation rule that validates client updates through hidden-task certificates built from server-side labeled reference data, replacing one-shot geometric trust signals with multi-task utility checks.
\item We formulate aggregation as a capped-simplex optimization problem that combines worst-task utility across hidden tasks, coordinate-rank anomaly suppression, and temporal regularization, yielding a robust weighting mechanism that remains compatible with heterogeneous FL.
\item We provide repository-backed empirical evidence for CARAT and clearly distinguish completed verified runs from the broader clean-accuracy and ASR benchmark suite that is still being finalized.
\end{itemize}

\section{Related Work}
\label{sec:related-work}

\paragraph{Model poisoning attacks.}
Modern FL attacks are explicitly designed to remain effective under server-side robustness measures. Early work already showed that highly constrained adversaries can poison the global model while preserving substantial stealth, even when the server deploys Byzantine-style defenses \citep{bhagoji2019adversarial,fang2020local,shejwalkar2021manipulating}. On the untargeted side, MinMax, MinSum, ALIE, and related optimization-based attacks craft malicious updates that evade simple geometric or statistical filters while still degrading global learning \citep{shejwalkar2021manipulating,baruch2019alie}. On the targeted side, model-replacement backdoors, distributed backdoor construction, trigger-tail attacks, and durable attacks such as Neurotoxin show that an update can appear globally plausible while implanting highly specific malicious behavior \citep{bagdasaryan2020backdoor,xie2020dba,wang2020tails,zhang2022neurotoxin}. Recent surveys make the broader point explicit: once the server cannot inspect raw client data, poisoning and backdoor attacks become a first-class systems concern rather than an edge case \citep{kairouz2021federated,liu2022threats}.

\paragraph{Defense and detection.}
Classical Byzantine-robust aggregation rules infer benignness from centrality. Krum and related geometry-based methods favor updates that remain close to other clients under pairwise distances \citep{blanchard2017krum,elmhamdi2018bulyan}, while coordinate-wise median and trimmed mean suppress marginal extremes independently \citep{yin2018median}. Later variants extend this family through suspicion scoring, regularized robust aggregation, coding-based redundancy, and geometric-median aggregation \citep{xie2019zeno,li2019rsa,chen2018draco,pillutla2022rfa,yin2019saddle}. These methods are appealing because they require limited server-side prior knowledge, but under heterogeneous data they may still confuse legitimate benign dispersion with malicious deviation. This limitation has motivated a second line of work that explicitly tackles heterogeneity, for example through FedProx- and SCAFFOLD-style stabilization, resampling, bucketing, local-SGD analysis, gradient splitting, nearest-neighbor mixing, label-skew-aware aggregation, and entropy-based public-data guidance \citep{li2020fedprox,karimireddy2020scaffold,he2021resampling,karimireddy2022bucketing,data2021localsgd,liu2023gas,allouah2023fixing,bao2024boba,huang2024sdea,zhu2021noniidsurvey}. A third line augments aggregation with trusted data, history, or representation-level inspection. In collaborative-learning settings, AUROR already anticipated update screening for poisoning resistance \citep{shen2016auror}; in FL, FLTrust introduces a trusted server anchor, FoolsGold and FLDetector rely on historical similarity diagnostics, and FLARE, FLAME, DeepSight, CRFL, and FLIP use latent-space inspection, clustering, clipping, smoothing, or reverse-engineering-based backdoor mitigation \citep{cao2021fltrust,fung2020foolsgold,zhang2022fldetector,wang2022flare,nguyen2022flame,rieger2022deepsight,xie2021crfl,zhang2023flip}. Recent work also revisits the cost of robustness itself, highlighting clipping design and the tradeoff between Byzantine resilience and no-attack accuracy \citep{allouah2025arc,yang2025tension}. CARAT is closest to this broader defense-and-detection family, but differs in its validation principle: instead of trusting update similarity, one reference direction, or post-hoc clustering alone, it evaluates whether a client update remains useful across multiple hidden semantic sub-tasks.

\section{Method}
\label{sec:method}

\subsection{Formulation and Assumptions}

We consider synchronous FL with $n$ clients and a global parameter vector $w_r \in \R^d$ at round $r$. At each round, client $i$ returns a message $u_i^r$, which the server maps to a vector-form update $\Delta_i^r \in \R^d$. For FedAvg, $\Delta_i^r$ can be read as the client model delta relative to the current global model; for gradient-style optimizers, it is the client update itself. The server must aggregate $\{\Delta_i^r\}_{i=1}^n$ into a single global step while a fraction $\rho$ of clients may be malicious and may coordinate their behavior.

CARAT assumes the server has access to labeled reference data $\mathcal{D}^{\mathrm{ref}}$. From this data, it constructs a probe pool $\mathcal{P}$ and repeatedly samples hidden, class-balanced validation tasks. This is a stronger assumption than fully data-free robust aggregation, and it should be treated as part of the method definition rather than as an incidental implementation detail. The reason for making this assumption explicit is that CARAT uses the server's auxiliary data not to define one trusted direction, but to generate many hidden semantic checks against which candidate client updates can be evaluated.

The goal is to assign aggregation weights that preserve benign diversity while suppressing updates that are either semantically harmful or structurally anomalous. This distinction matters in non-i.i.d.\ FL: benign clients may differ substantially in direction, norm, and class emphasis, so geometric dispersion alone is not enough to identify attacks. CARAT therefore separates each round into three conceptual steps. It first places all client updates into a common evaluation geometry, then measures each client's task utility and structural abnormality, and finally converts these signals into robust aggregation weights. This separation is central to the method: evaluation should be comparable across clients, while the final server step should still remain conservative and robust.

\subsection{Method Details}

At round $r$, CARAT begins by converting all client messages into vector-form deltas and computing their norms. Before any semantic evaluation, the server robustly clips client updates to prevent a small number of large-magnitude messages from dominating the comparison stage. Let $\|\Delta_i\|_2$ denote the norm of client $i$'s update. In our implementation, the clipping threshold $\tau_r$ is estimated by a MAD-based rule:
\begin{equation}
\tau_r = \operatorname{median}(\{\|\Delta_i\|_2\}) +
k \cdot 1.4826 \cdot \operatorname{MAD}(\{\|\Delta_i\|_2\}),
\label{eq:clip}
\end{equation}
and the clipped update is
\begin{equation}
\bar{\Delta}_i = \min\left(1, \frac{\tau_r}{\|\Delta_i\|_2 + \varepsilon}\right)\Delta_i.
\end{equation}

CARAT then rescales each clipped update to a common evaluation radius $q_r$:
\begin{equation}
\tilde{\Delta}_i = q_r \frac{\bar{\Delta}_i}{\|\bar{\Delta}_i\|_2 + \varepsilon},
\label{eq:normalize}
\end{equation}
where $q_r$ is chosen as a quantile of the positive clipped norms. This normalization removes pure scale advantages during hidden-task comparison. Importantly, the normalized vectors are used only for evaluation; the final model update is still computed from the clipped vectors. In other words, CARAT uses normalization to create a fair scoring geometry, not to define the actual optimization step taken by the server.

After this preprocessing stage, CARAT samples $T$ hidden tasks from the probe pool $\mathcal{P}$, with each task constructed to be approximately class-balanced. This is the source of the method's class-aware behavior. A malicious update can remain globally plausible while degrading only a subset of classes or semantic regions, whereas repeated balanced hidden tasks make these localized failures harder to hide. For each hidden task $t$, let $\ell_t(w)$ denote the average cross-entropy loss of model $w$ on that task. CARAT defines a \emph{hidden-task certificate}
\begin{equation}
c_{i,t} = \ell_t(w_r) - \ell_t(w_r + \tilde{\Delta}_i).
\label{eq:cert}
\end{equation}
Here, ``certificate'' denotes an internal task-level utility score rather than a formal certified-robustness guarantee. Positive values indicate that client $i$ improves task $t$ relative to the current global model.

Stacking these scores yields a certificate matrix $C \in \R^{n \times T}$. Because different tasks may have different intrinsic difficulty and baseline loss scale, CARAT robustly standardizes each column of $C$ before optimization. This preserves within-task ordering while making the task scores comparable across different hidden slices. The resulting matrix gives the server a much richer signal than client-to-client similarity alone: a client should be trusted because it remains broadly useful across hidden tasks, not merely because it lies near the geometric center of the client population.

Task utility alone is still insufficient, because a malicious client may overfit to the probe tasks or align with part of the semantic signal while remaining structurally extreme in parameter space. CARAT therefore augments hidden-task certificates with a coordinate-rank anomaly score. For a sampled coordinate set $S \subseteq \{1,\dots,d\}$, it ranks clients coordinate-wise and measures the average normalized distance from the center rank:
\begin{equation}
r_i^{(S)} = \frac{1}{|S|}\sum_{j \in S}
\left|
\frac{\operatorname{rank}_{i,j}}{n-1} - \frac{1}{2}
\right|.
\label{eq:rank}
\end{equation}
CARAT repeats this computation over several random coordinate subsets and aggregates the resulting values robustly. A client receives a large anomaly score if its update remains persistently extreme across many coordinates rather than only on a few isolated entries. After robust standardization, only positive deviations are retained as penalties. The signal is intentionally asymmetric: suspicious extremeness is penalized, whereas unusual centrality is not explicitly rewarded. Together, the certificate matrix and rank penalty separate two notions that are often conflated in robust aggregation: whether an update is useful, and whether it remains structurally atypical.

\subsection{Weight Optimization and Aggregation}

CARAT converts these signals into aggregation weights by solving a constrained optimization problem on the capped simplex
\begin{equation}
\mathcal{A}(\hat{\rho}) =
\left\{
\alpha \in \R^n :
\alpha_i \ge 0,\;
\sum_{i=1}^n \alpha_i = 1,\;
\alpha_i \le \frac{1}{(1-\hat{\rho})n}
\right\},
\label{eq:capped-simplex}
\end{equation}
where $\hat{\rho}$ is an estimate of the attacker fraction. The cap prevents any single client from dominating the aggregate and encodes a conservative trust budget at the weight level: even when a client appears useful, it cannot absorb arbitrary influence in one round.

Let $r \in \R^n$ denote the positive rank penalties and let $\alpha^{\mathrm{prev}}$ be the previous round's weights. CARAT solves
\begin{equation}
\max_{\alpha \in \mathcal{A}(\hat{\rho})}
\;\softmin_{\beta}(C^\top \alpha)
- \lambda_r r^\top \alpha
- \lambda_p \|\alpha - \alpha^{\mathrm{prev}}\|_2^2,
\label{eq:objective}
\end{equation}
where
\begin{equation}
\softmin_{\beta}(s_1,\dots,s_T)
= -\frac{1}{\beta}\log\left(\frac{1}{T}\sum_{t=1}^T e^{-\beta s_t}\right).
\label{eq:softmin}
\end{equation}
The first term is a smooth approximation to worst-task utility: large $\beta$ places more emphasis on the weakest hidden tasks, while finite $\beta$ preserves optimization stability. The second term penalizes structurally extreme clients, and the third term regularizes the solution toward the previous round to reduce unnecessary oscillation. Eq.~\eqref{eq:objective} therefore encodes a clear decision rule: a client should retain influence only if it remains useful across hidden tasks and does not repeatedly appear as an outlier in rank space.

After optimization, CARAT aggregates the clipped updates:
\begin{equation}
\Delta_{\mathrm{agg}} = \sum_{i=1}^n \alpha_i \bar{\Delta}_i.
\label{eq:aggregate}
\end{equation}
This yields the final server step. The final update therefore inherits the robustness benefits of clipping while using weights learned from the richer evaluation geometry defined above. In practice, the main additional computation comes from evaluating $n$ candidate updates over $T$ hidden tasks and computing rank penalties on coordinate subsamples; the probe-model forward passes dominate, while the ranking and optimization stages remain tunable through explicit budget parameters.

A compact round-level algorithm summary is provided in Appendix~\ref{sec:appendix-algorithm}. Formal properties of the normalization step, capped-simplex feasibility, softmin behavior, temporal regularization, and per-round complexity are deferred to Appendix~\ref{sec:appendix-theory}.

\section{Experiments}
\label{sec:experiments}

\subsection{Experimental Setup}

\paragraph{Datasets and Implementation.}
Our evaluation uses two widely adopted benchmark datasets:  CIFAR-10 \cite{krizhevsky2009learning} and CIFAR-100 \cite{krizhevsky2009learning}. To emulate a realistic federated environment with heterogeneous data, we partition samples among clients using a Dirichlet distribution ( $\bm{p}\sim \text{Dir}_N(\alpha)$ ). We set the concentration parameter to $\alpha=0.1, 0.5, 1$,  standard choices in the literature that induces a significant degree of non-IID distribution in the local datasets of the clients.

The experiments employ two representative neural network architectures to assess performance across varying model complexities: VGG19~\cite{simonyan2015deepconvolutionalnetworkslargescale}, and ResNet18~\cite{He_2016_CVPR}. Specifically, we pair the VGG19 model with the CIFAR-10 dataset. For the more complex CIFAR-100 dataset, we evaluate performance using ResNet18 to cover a diverse set of architectures.

\paragraph{Baselines and Default Setting.}
To benchmark the performance of \texttt{TriGuardFL}, we compare it against four methods:  FLTrust~\cite{cao2021fltrust}, FLDetector~\cite{zhang2022fldetector}, NormClipping~\cite{sun2019normclipping}, and MultiKrum ~\cite{blanchard2017machine}. With regard to the attack models, four model poisoning attacks are adopted: Min-Sum~\cite{shejwalkar2021manipulating}, Min-Max \cite{shejwalkar2021manipulating}, ALIE \cite{baruch2019little}, and FangAttack \cite{fang2020local}.

The clean-reference benchmark uses \texttt{NoAttack} and compares \{Mean, FLTrust, MultiKrum, CARAT\}. The untargeted poisoning benchmark evaluates \{ALIE, MinMax, FangAttack\} against the same defense set. The backdoor benchmark evaluates \{BadNets, Neurotoxin\} against \{Mean, FLTrust, FLAME, DeepSight, CARAT\}. We consider both i.i.d.\ partitions and Dirichlet non-i.i.d.\ splits: \texttt{clean\_reference.yaml} and \texttt{main\_untargeted.yaml} use $\alpha \in \{0.1,0.5\}$, while \texttt{backdoor.yaml} uses i.i.d.\ plus $\alpha=0.5$. The untargeted suite sweeps attacker fractions $\{0.2,0.4\}$; the clean-reference and backdoor suites use $0$ and $0.2$, respectively. Each paper-side suite runs three repeats with effective seeds $42$, $43$, and $44$. The default CARAT hyperparameters are: probe pool size $512$, $8$ hidden tasks, $4$ probe samples per class, probe batch size $128$, probe-radius quantile $0.5$, MAD clipping with $k=2.5$, default attacker prior $\hat{\rho}=0.2$, certificate temperature $\beta=12.0$, rank weight $0.20$, prior weight $0.05$, $4096$ sampled coordinates, $5$ rank subsamples, $60$ optimization steps, and optimizer learning rate $0.5$. One important caveat is that the strongest completed legacy CARAT pilot was produced from an earlier base config with \verb|evaluate=false|. Its main-text evidence is therefore limited to training accuracy, training loss, and runtime, whereas all newly committed benchmark suites are configured to record test-time metrics.

The key hyperparameters used throughout the experiments are summarized in Table \ref{tab:hyperparams}.
\begin{table*}[!htbp]
\centering
\caption{Default hyperparameter configuration.}
\label{tab:hyperparams}
\renewcommand{\arraystretch}{1}
    \setlength{\tabcolsep}{3pt}
\small
\begin{tabular}{cccccccccc}
\toprule
\textbf{Parameter} &
\makecell[c]{\textbf{Total}\\\textbf{Clients} ($N$)} &
\makecell[c]{\textbf{Malicious}\\\textbf{Clients Ratio} ($\%$)} &
\makecell[c]{\textbf{Clients}\\\textbf{per Round}} &
\makecell[c]{\textbf{Local}\\\textbf{Iterations} ($K$)} &
\makecell[c]{\textbf{Total}\\\textbf{Epochs} ($T$)} &
\makecell[c]{\textbf{Learning}\\\textbf{Rate} ($\eta$)} &
\makecell[c]{\textbf{Significance}\\\textbf{Level}}  \\
\midrule
\textbf{Value} & 20 & 10/20/30  & 20 & 4 & 200& 0.05 & 0.02 \\
\bottomrule
\end{tabular}
\end{table*}

\subsection{Performance Evaluation}

\paragraph{Performance of Defense Effectiveness under Various Attacks.}
\subsection{Mitigating the impact of attacks} 
  \begin{table*}[!htbp]
      \definecolor{mygray}{gray}{.90}
      \renewcommand{\arraystretch}{1.0}
      \setlength{\tabcolsep}{4pt}
      \caption{Loss under four state-of-the-art model poisoning attacks with $10\%$
  malicious clients on two datasets and model architectures under non-IID settings.}
      \label{tablossmergedcifar10cifar100}
      \centering\small
      \begin{tabular}{c c c c c c c c}
      \toprule
      \multirow{2}{*}{\makecell[c]{Dataset\\Model}} &
      \multirow{2}{*}{\makecell[c]{Non-IID\\Settings}} &
      \multirow{2}{*}{\makecell[c]{Defense\\Rules}} &
      \multicolumn{4}{c}{Attacks} &
      \multirow{2}{*}{Rank} \\
      \cmidrule(lr){4-7}
      & & & ALIE & FangAttack & MinMax & MinSum & \\
      \midrule

      \multirow{18}{*}{\makecell[c]{CIFAR-10\\VGG19}}
      & \multirow{6}{*}{\makecell[c]{$\alpha=1$}}
      & NormClipping & 0.0176$\pm$0.0003 & 0.0184$\pm$0.0003 & 0.0176$\pm$0.0003 & 0.0177$
  \pm$0.0003 & 5.0 \\
      & & MultiKrum & 0.0120$\pm$0.0001 & 0.0159$\pm$0.0009 & 0.0157$\pm$0.0005 & 0.0153$
  \pm$0.0010 & 4.0 \\
      & & FLTrust & 0.0152$\pm$0.0003 & 0.0152$\pm$0.0003 & 0.0148$\pm$0.0003 & 0.0149$
  \pm$0.0003 & 3.0 \\
      & & FLDetector & \textbf{0.0093$\pm$0.0002} & 0.0192$\pm$0.0018 & \textbf{0.0089$
  \pm$0.0001} & \textbf{0.0089$\pm$0.0002} & 2.0 \\
      & & \cellcolor{mygray}TriGuardFL & \cellcolor{mygray}\textbf{0.0093$\pm$0.0002} &
  \cellcolor{mygray}\textbf{0.0093$\pm$0.0002} & \cellcolor{mygray}0.0090$\pm$0.0002 &
  \cellcolor{mygray}0.0090$\pm$0.0002 & \cellcolor{mygray}1.0 \\
      & & NoAttack+Mean(Base) & 0.0092$\pm$0.0001 & 0.0092$\pm$0.0001 & 0.0092$\pm$0.0001 &
  0.0092$\pm$0.0001 & -- \\
      \cmidrule(lr){2-8}

      & \multirow{6}{*}{\makecell[c]{$\alpha=0.5$}}
      & NormClipping & 0.0194$\pm$0.0002 & 0.0200$\pm$0.0004 & 0.0196$\pm$0.0004 & 0.0195$
  \pm$0.0003 & 4.5 \\
      & & MultiKrum & 0.0147$\pm$0.0008 & 0.0207$\pm$0.0021 & 0.0196$\pm$0.0012 & 0.0208$
  \pm$0.0023 & 4.5 \\
      & & FLTrust & 0.0167$\pm$0.0004 & 0.0168$\pm$0.0003 & 0.0167$\pm$0.0004 & 0.0163$
  \pm$0.0003 & 3.0 \\
      & & FLDetector & 0.0108$\pm$0.0005 & 0.0242$\pm$0.0005 & 0.0110$\pm$0.0010 & 0.0106$
  \pm$0.0006 & 2.0 \\
      & & \cellcolor{mygray}TriGuardFL & \cellcolor{mygray}\textbf{0.0106$\pm$0.0003} &
  \cellcolor{mygray}\textbf{0.0107$\pm$0.0003} & \cellcolor{mygray}\textbf{0.0107$
  \pm$0.0010} & \cellcolor{mygray}\textbf{0.0104$\pm$0.0003} & \cellcolor{mygray}1.0 \\
      & & NoAttack+Mean(Base) & 0.0104$\pm$0.0004 & 0.0104$\pm$0.0004 & 0.0104$\pm$0.0004 &
  0.0104$\pm$0.0004 & -- \\
      \cmidrule(lr){2-8}

      & \multirow{6}{*}{\makecell[c]{$\alpha=0.1$}}
      & NormClipping & 0.0256$\pm$0.0011 & 0.0261$\pm$0.0012 & 0.0268$\pm$0.0013 & 0.0260$
  \pm$0.0015 & 4.0 \\
      & & MultiKrum & 0.0270$\pm$0.0013 & 0.0330$\pm$0.0028 & 0.0321$\pm$0.0034 & 0.0327$
  \pm$0.0026 & 5.0 \\
      & & FLTrust & 0.0233$\pm$0.0011 & 0.0236$\pm$0.0013 & 0.0243$\pm$0.0018 & 0.0239$
  \pm$0.0013 & 3.0 \\
      & & FLDetector & 0.0215$\pm$0.0022 & 0.0243$\pm$0.0022 & 0.0218$\pm$0.0031 & 0.0218$
  \pm$0.0034 & 2.0 \\
      & & \cellcolor{mygray}TriGuardFL & \cellcolor{mygray}\textbf{0.0198$\pm$0.0016} &
  \cellcolor{mygray}\textbf{0.0198$\pm$0.0018} & \cellcolor{mygray}\textbf{0.0198$
  \pm$0.0016} & \cellcolor{mygray}\textbf{0.0197$\pm$0.0024} & \cellcolor{mygray}1.0 \\
      & & NoAttack+Mean(Base) & 0.0199$\pm$0.0012 & 0.0199$\pm$0.0012 & 0.0199$\pm$0.0012 &
  0.0199$\pm$0.0012 & -- \\
      \midrule

      \multirow{18}{*}{\makecell[c]{CIFAR-100\\ResNet18}}
      & \multirow{6}{*}{\makecell[c]{$\alpha=1$}}
      & NormClipping & 0.0473$\pm$0.0003 & 0.0487$\pm$0.0004 & 0.0479$\pm$0.0003 & 0.0480$
  \pm$0.0003 & 5.0 \\
      & & MultiKrum & 0.0330$\pm$0.0004 & 0.0374$\pm$0.0008 & 0.0375$\pm$0.0005 & 0.0372$
  \pm$0.0007 & 3.0 \\
      & & FLTrust & 0.0405$\pm$0.0003 & 0.0405$\pm$0.0004 & 0.0403$\pm$0.0005 & 0.0404$
  \pm$0.0004 & 4.0 \\
      & & FLDetector & \textbf{0.0274$\pm$0.0005} & 0.0426$\pm$0.0004 & 0.0275$\pm$0.0004 &
  \textbf{0.0275$\pm$0.0005} & 2.0 \\
      & & \cellcolor{mygray}TriGuardFL & \cellcolor{mygray}0.0276$\pm$0.0004 &
  \cellcolor{mygray}\textbf{0.0276$\pm$0.0004} & \cellcolor{mygray}\textbf{0.0273$
  \pm$0.0003} & \cellcolor{mygray}\textbf{0.0275$\pm$0.0002} & \cellcolor{mygray}1.0 \\
      & & NoAttack+Mean(Base) & 0.0269$\pm$0.0003 & 0.0269$\pm$0.0003 & 0.0269$\pm$0.0003 &
  0.0269$\pm$0.0003 & -- \\
      \cmidrule(lr){2-8}

      & \multirow{6}{*}{\makecell[c]{$\alpha=0.5$}}
      & NormClipping & 0.0482$\pm$0.0006 & 0.0496$\pm$0.0003 & 0.0488$\pm$0.0004 & 0.0490$
  \pm$0.0005 & 5.0 \\
      & & MultiKrum & 0.0352$\pm$0.0008 & 0.0403$\pm$0.0015 & 0.0405$\pm$0.0012 & 0.0402$
  \pm$0.0010 & 3.0 \\
      & & FLTrust & 0.0414$\pm$0.0005 & 0.0414$\pm$0.0004 & 0.0412$\pm$0.0004 & 0.0413$
  \pm$0.0005 & 4.0 \\
      & & FLDetector & \textbf{0.0284$\pm$0.0004} & 0.0437$\pm$0.0011 & \textbf{0.0284$
  \pm$0.0004} & \textbf{0.0283$\pm$0.0004} & 2.0 \\
      & & \cellcolor{mygray}TriGuardFL & \cellcolor{mygray}\textbf{0.0284$\pm$0.0004} &
  \cellcolor{mygray}\textbf{0.0285$\pm$0.0005} & \cellcolor{mygray}\textbf{0.0284$
  \pm$0.0006} & \cellcolor{mygray}\textbf{0.0283$\pm$0.0005} & \cellcolor{mygray}1.0 \\
      & & NoAttack+Mean(Base) & 0.0276$\pm$0.0005 & 0.0276$\pm$0.0005 & 0.0276$\pm$0.0005 &
  0.0276$\pm$0.0005 & -- \\
      \cmidrule(lr){2-8}

      & \multirow{6}{*}{\makecell[c]{$\alpha=0.1$}}
      & NormClipping & 0.0519$\pm$0.0004 & 0.0529$\pm$0.0004 & 0.0527$\pm$0.0005 & 0.0529$
  \pm$0.0004 & 5.0 \\
      & & MultiKrum & 0.0424$\pm$0.0010 & 0.0504$\pm$0.0010 & 0.0504$\pm$0.0015 & 0.0504$
  \pm$0.0014 & 4.0 \\
      & & FLTrust & 0.0450$\pm$0.0004 & 0.0450$\pm$0.0005 & 0.0450$\pm$0.0005 & 0.0451$
  \pm$0.0006 & 3.0 \\
      & & FLDetector & 0.0326$\pm$0.0003 & 0.0461$\pm$0.0039 & 0.0325$\pm$0.0003 &
  \textbf{0.0323$\pm$0.0005} & 2.0 \\
      & & \cellcolor{mygray}TriGuardFL & \cellcolor{mygray}\textbf{0.0322$\pm$0.0004} &
  \cellcolor{mygray}\textbf{0.0323$\pm$0.0004} & \cellcolor{mygray}\textbf{0.0322$
  \pm$0.0005} & \cellcolor{mygray}\textbf{0.0323$\pm$0.0005} & \cellcolor{mygray}1.0 \\
      & & NoAttack+Mean(Base) & 0.0314$\pm$0.0004 & 0.0314$\pm$0.0004 & 0.0314$\pm$0.0004 &
  0.0314$\pm$0.0004 & -- \\
      \bottomrule
      \end{tabular}
  \end{table*}


We subjected the FL system to four potent attacks: Min-Max, Min-Sum, ALIE, and FangAttack.

Table~\ref{tablossmergedcifar10cifar100} presents the final test loss for all evaluated defenses across the different network architectures. According to the average performance rank, \texttt{TriGuardFL} is the highest-performing defense, achieving rank 1 in three different non-IID experimental settings. This resilience is particularly evident under strong and adaptive attacks. For example, when defending a ResNet-18 model against the FangAttack, \texttt{TriGuardFL} maintained a test loss of $0.0323$; meanwhile, the losses for MultiKrum and FLTrust significantly degraded to $0.0504$ and $0.0450$, respectively. 

\begin{figure*}[!htbp]
	\centering
	% \subfloat[IID data distribution]{\includegraphics[trim=0 0 0 2.5cm,
 %  clip,scale=0.31]{figures/iid_CIFAR100_resnet18_FedAvg_adv0.2_alphaNA_clients20_lr0.05_cfgFedAvg_CIFAR100_config_epochs200.pdf}}
	% \quad
	\subfloat[Non-IID data distribution, $\alpha=1$]{\includegraphics[trim=0 0 0 2.5cm,
  clip,scale=0.31]{figures/non-iid_CIFAR100_resnet18_FedAvg_adv0.2_alpha1_clients20_lr0.05_cfgFedAvg_CIFAR100_config_epochs200.pdf}}
	\quad
	\subfloat[Non-IID data distribution, $\alpha=0.5$]{\includegraphics[trim=0 0 0 2.5cm,
  clip,scale=0.31]{figures/non-iid_CIFAR100_resnet18_FedAvg_adv0.2_alpha0.5_clients20_lr0.05_cfgFedAvg_CIFAR100_config_epochs200.pdf}}
    \quad
	\subfloat[Non-IID data distribution, $\alpha=0.1$]{\includegraphics[trim=0 0 0 2.5cm,
  clip,scale=0.31]{figures/non-iid_CIFAR100_resnet18_FedAvg_adv0.2_alpha0.1_clients20_lr0.05_cfgFedAvg_CIFAR100_config_epochs200.pdf}}
    \caption{Evolution of test accuracy on CIFAR-100, comparing \texttt{TriGuardFL} against baseline defenses under four state-of-the-art  model poisoning attacks. The experimental setup uses a ResNet-18 model trained among $20$ total clients, including $4$ malicious attackers.}


	\label{figaccuracycifar100}
\end{figure*}



Fig.~\ref{figaccuracycifar100} shows that the accuracy curve for \texttt{TriGuardFL} remains high and stable throughout the training process, closely tracking the performance of the ideal ``no attack" baseline. This stands in contrast to several baseline defenses, such as FLTrust and MultiKrum, which exhibit significant performance degradation and instability, particularly under the more sophisticated full-knowledge attacks. These results underscore \texttt{TriGuardFL}'s ability to effectively filter out malicious updates without incorrectly penalizing benign clients, which is a critical weakness in many defenses that struggle with non-IID data.

\paragraph{Impact of Attacker Ratio.}
The proportion of malicious clients significantly influences the accuracy of the FL model accuracy. In our previous experiments, we set $10\%$ ratio for malicious clients as default. To evaluate the detection ability and model accuracy of \texttt{TriGuardFL} under different numbers of adversaries, we conducted additional tests with $20\%$ and $30\%$ proportion for malicious clients, while keeping other settings constant. As shown in %Table \ref{tabmoreattackers} and 
Fig. \ref{figaccuracymoreattackerscifar100}, the accuracy of the test dataset decreases %and the losses of the test dataset increase 
in general. For \texttt{TriGuardFL}, the losses of the test dataset do not have significant differences with different numbers of malicious clients, indicating that \texttt{TriGuardFL} keeps excellent robustness. Moreover, FLTrust and MultiKrum crash in part of the cases. This is because even if \texttt{TriGuardFL} can detect most malicious clients, the number of benign clients decreases as that of malicious clients increases, and samples in benign clients also decrease. 

\begin{figure*}[htbp]
	\centering
    % \subfloat[Dataset distribution is IID.]{\includegraphics[trim=0 0 0 2.5cm,clip,scale=0.31]{figures/adv_sensitivity_iid_CIFAR100_resnet18_FedAvg_alphaNA_clients20_lr0.05_cfgFedAvg_CIFAR100_config_advset0.1-0.2-0.3_epochs200.pdf}}
  
	% \quad
	\subfloat[Dataset distribution is non-IID with $\alpha=1$.]{\includegraphics[trim=0 0 0 2.5cm,
  clip,scale=0.31]{figures/adv_sensitivity_non-iid_CIFAR100_resnet18_FedAvg_alpha1_clients20_lr0.05_cfgFedAvg_CIFAR100_config_advset0.1-0.2-0.3_epochs200.pdf}}
	\quad
	\subfloat[Dataset distribution is non-IID with $\alpha=0.5$.]{\includegraphics[trim=0 0 0 2.5cm,
  clip,scale=0.31]{figures/adv_sensitivity_non-iid_CIFAR100_resnet18_FedAvg_alpha0.5_clients20_lr0.05_cfgFedAvg_CIFAR100_config_advset0.1-0.2-0.3_epochs200.pdf}}
  \quad
	\subfloat[Dataset distribution is non-IID with $\alpha=0.1$.]{\includegraphics[trim=0 0 0 2.5cm,
  clip,scale=0.31]{figures/adv_sensitivity_non-iid_CIFAR100_resnet18_FedAvg_alpha0.1_clients20_lr0.05_cfgFedAvg_CIFAR100_config_advset0.1-0.2-0.3_epochs200.pdf}}
	\caption{Accuracy sensitivity ratio of different adversary settings comparing \texttt{TriGuardFL} and another four defense algorithms when the test dataset is CIFAR-100 and the neural network model is ResNet18.}
	\label{figaccuracymoreattackerscifar100}
\end{figure*}

\paragraph{Impact of the number of clients-scale}
\begin{figure*}[!htbp]
	\centering
	\subfloat[Dataset distribution is non-IID with $\alpha=1$.]{\includegraphics[trim=0 0 0 2.5cm,
  clip,scale=0.31]{figures/num_clients_sensitivity_non-iid_CIFAR100_resnet18_FedAvg_alpha1_adv0.2_lr0.05_cfgFedAvg_CIFAR100_config_clientset20-40-60_epochs200.pdf}}
	\quad
	\subfloat[Dataset distribution is non-IID with $\alpha=0.5$.]{\includegraphics[trim=0 0 0 2.5cm,
  clip,scale=0.31]{figures/num_clients_sensitivity_non-iid_CIFAR100_resnet18_FedAvg_alpha0.5_adv0.2_lr0.05_cfgFedAvg_CIFAR100_config_clientset20-40-60_epochs200.pdf}}
  \quad
	\subfloat[Dataset distribution is non-IID with $\alpha=0.1$.]{\includegraphics[trim=0 0 0 2.5cm,
  clip,scale=0.31]{figures/num_clients_sensitivity_non-iid_CIFAR100_resnet18_FedAvg_alpha0.1_adv0.2_lr0.05_cfgFedAvg_CIFAR100_config_clientset20-40-60_epochs200.pdf}}
	\caption{Accuracy sensitivity of different clients settings comparing \texttt{TriGuardFL} and another four defense algorithms when the test dataset is CIFAR-100 and the neural network model is ResNet18.}
	\label{figaccuracymoreclientscifar100}
\end{figure*}
  \begin{table*}[!h]
      \definecolor{mygray}{gray}{.90}
      \renewcommand{\arraystretch}{1.0}
      \setlength{\tabcolsep}{4pt}
      \caption{Loss of the CIFAR-100 dataset with ResNet18 under four state-of-the-art model
  poisoning attacks at adversary ratio $20\%$, comparing non-IID settings across different
  client counts. Each cell reports mean$\pm$std over available seeds. If a run stopped
  early, its last recorded epoch before 200 is used.}
      \label{tablosscifar100clientsadv02}
      \centering\small
      \begin{tabular}{c c c c c c c c}
      \toprule
      \multirow{2}{*}{\makecell[c]{Number of\\Clients}} &
      \multirow{2}{*}{\makecell[c]{Non-IID\\Settings}} &
      \multirow{2}{*}{\makecell[c]{Defense\\Rules}} &
      \multicolumn{4}{c}{Attacks} &
      \multirow{2}{*}{\makecell[c]{Rank}} \\
      \cmidrule(lr){4-7}\cmidrule(lr){8-8}
      & & & ALIE & FangAttack & MinMax & MinSum & \\
      \midrule

      \multirow{18}{*}{40}
      & \multirow{6}{*}{\makecell[c]{$\alpha=1$}}
      & NormClipping & 0.0472$\pm$0.0003 & 0.0586$\pm$0.0067 & 0.0479$\pm$0.0002 & 0.0484$
  \pm$0.0003 & 5.0 \\
      & & MultiKrum & 0.0427$\pm$0.0012 & 0.0309$\pm$0.0006 & 0.0268$\pm$0.0007 & 0.0284$
  \pm$0.0013 & 3.0 \\
      & & FLTrust & 0.0412$\pm$0.0002 & 0.0405$\pm$0.0001 & 0.0403$\pm$0.0001 & 0.0405$
  \pm$0.0003 & 4.0 \\
      & & FLDetector & 0.0299$\pm$0.0007 & 0.0601$\pm$0.0024 & 0.0259$\pm$0.0005 & 0.0260$
  \pm$0.0006 & 2.0 \\
      & & \cellcolor{mygray}TriGuardFL & \cellcolor{mygray}\textbf{0.0276$\pm$0.0006} &
      \cellcolor{mygray}\textbf{0.0258$\pm$0.0006} &
      \cellcolor{mygray}\textbf{0.0257$\pm$0.0007} &
      \cellcolor{mygray}\textbf{0.0257$\pm$0.0005} & \cellcolor{mygray}1.0 \\
      & & NoAttack+Mean(Base) & 0.0251$\pm$0.0006 & 0.0251$\pm$0.0006 & 0.0251$\pm$0.0006 &
  0.0251$\pm$0.0006 & -- \\
      \cmidrule(lr){2-8}

      & \multirow{6}{*}{\makecell[c]{$\alpha=0.5$}}
      & NormClipping & 0.0482$\pm$0.0005 & 0.0560$\pm$0.0005 & 0.0490$\pm$0.0006 & 0.0496$
  \pm$0.0008 & 5.0 \\
      & & MultiKrum & 0.0448$\pm$0.0008 & 0.0327$\pm$0.0005 & 0.0279$\pm$0.0005 & 0.0293$
  \pm$0.0013 & 3.0 \\
      & & FLTrust & 0.0422$\pm$0.0002 & 0.0415$\pm$0.0004 & 0.0414$\pm$0.0005 & 0.0416$
  \pm$0.0004 & 4.0 \\
      & & FLDetector & 0.0311$\pm$0.0006 & 0.0572$\pm$0.0032 & 0.0267$\pm$0.0003 & 0.0268$
  \pm$0.0004 & 2.0 \\
      & & \cellcolor{mygray}TriGuardFL & \cellcolor{mygray}\textbf{0.0290$\pm$0.0004} &
      \cellcolor{mygray}\textbf{0.0267$\pm$0.0004} &
      \cellcolor{mygray}\textbf{0.0266$\pm$0.0003} &
      \cellcolor{mygray}\textbf{0.0265$\pm$0.0002} & \cellcolor{mygray}1.0 \\
      & & NoAttack+Mean(Base) & 0.0259$\pm$0.0003 & 0.0259$\pm$0.0003 & 0.0259$\pm$0.0003 &
  0.0259$\pm$0.0003 & -- \\
      \cmidrule(lr){2-8}

      & \multirow{6}{*}{\makecell[c]{$\alpha=0.1$}}
      & NormClipping & 0.0524$\pm$0.0004 & 0.0602$\pm$0.0005 & 0.0537$\pm$0.0006 & 0.0544$
  \pm$0.0006 & 5.0 \\
      & & MultiKrum & 0.0552$\pm$0.0007 & 0.0419$\pm$0.0010 & 0.0346$\pm$0.0005 & 0.0370$
  \pm$0.0014 & 3.0 \\
      & & FLTrust & 0.0466$\pm$0.0005 & 0.0458$\pm$0.0005 & 0.0460$\pm$0.0005 & 0.0468$
  \pm$0.0005 & 4.0 \\
      & & FLDetector & 0.0373$\pm$0.0009 & 0.0550$\pm$0.0043 & 0.0315$\pm$0.0008 & 0.0315$
  \pm$0.0006 & 2.0 \\
      & & \cellcolor{mygray}TriGuardFL & \cellcolor{mygray}\textbf{0.0345$\pm$0.0006} &
      \cellcolor{mygray}\textbf{0.0305$\pm$0.0004} &
      \cellcolor{mygray}\textbf{0.0311$\pm$0.0007} &
      \cellcolor{mygray}\textbf{0.0308$\pm$0.0005} & \cellcolor{mygray}1.0 \\
      & & NoAttack+Mean(Base) & 0.0303$\pm$0.0006 & 0.0303$\pm$0.0006 & 0.0303$\pm$0.0006 &
  0.0303$\pm$0.0006 & -- \\
      \midrule

      \multirow{18}{*}{60}
      & \multirow{6}{*}{\makecell[c]{$\alpha=1$}}
      & NormClipping & 0.0470$\pm$0.0003 & 0.0593$\pm$0.0037 & 0.0484$\pm$0.0012 & 0.0486$
  \pm$0.0003 & 5.0 \\
      & & MultiKrum & 0.0452$\pm$0.0003 & 0.0303$\pm$0.0011 & 0.0258$\pm$0.0005 & 0.0275$
  \pm$0.0009 & 3.0 \\
      & & FLTrust & 0.0410$\pm$0.0003 & 0.0427$\pm$0.0027 & 0.0401$\pm$0.0003 & 0.0405$
  \pm$0.0005 & 4.0 \\
      & & FLDetector & 0.0290$\pm$0.0003 & 0.0484$\pm$0.0007 & \textbf{0.0247$\pm$0.0003} &
  0.0249$\pm$0.0004 & 2.0 \\
      & & \cellcolor{mygray}TriGuardFL & \cellcolor{mygray}\textbf{0.0271$\pm$0.0003} &
      \cellcolor{mygray}\textbf{0.0260$\pm$0.0011} &
      \cellcolor{mygray}\textbf{0.0247$\pm$0.0003} &
      \cellcolor{mygray}\textbf{0.0247$\pm$0.0002} & \cellcolor{mygray}1.0 \\
      & & NoAttack+Mean(Base) & 0.0245$\pm$0.0004 & 0.0245$\pm$0.0004 & 0.0245$\pm$0.0004 &
  0.0245$\pm$0.0004 & -- \\
      \cmidrule(lr){2-8}

      & \multirow{6}{*}{\makecell[c]{$\alpha=0.5$}}
      & NormClipping & 0.0479$\pm$0.0004 & 0.0575$\pm$0.0011 & 0.0487$\pm$0.0003 & 0.0496$
  \pm$0.0004 & 5.0 \\
      & & MultiKrum & 0.0472$\pm$0.0008 & 0.0316$\pm$0.0006 & 0.0273$\pm$0.0013 & 0.0289$
  \pm$0.0013 & 3.0 \\
      & & FLTrust & 0.0422$\pm$0.0005 & 0.0431$\pm$0.0021 & 0.0411$\pm$0.0005 & 0.0418$
  \pm$0.0005 & 4.0 \\
      & & FLDetector & 0.0303$\pm$0.0006 & 0.0449$\pm$0.0012 & 0.0260$\pm$0.0005 & 0.0260$
  \pm$0.0004 & 2.0 \\
      & & \cellcolor{mygray}TriGuardFL & \cellcolor{mygray}\textbf{0.0285$\pm$0.0006} &
      \cellcolor{mygray}\textbf{0.0264$\pm$0.0004} &
      \cellcolor{mygray}\textbf{0.0255$\pm$0.0003} &
      \cellcolor{mygray}\textbf{0.0255$\pm$0.0003} & \cellcolor{mygray}1.0 \\
      & & NoAttack+Mean(Base) & 0.0253$\pm$0.0004 & 0.0253$\pm$0.0004 & 0.0253$\pm$0.0004 &
  0.0253$\pm$0.0004 & -- \\
      \cmidrule(lr){2-8}

      & \multirow{6}{*}{\makecell[c]{$\alpha=0.1$}}
      & NormClipping & 0.0518$\pm$0.0004 & 0.0641$\pm$0.0059 & 0.0530$\pm$0.0005 & 0.0538$
  \pm$0.0004 & 5.0 \\
      & & MultiKrum & 0.0563$\pm$0.0005 & 0.0389$\pm$0.0016 & 0.0327$\pm$0.0006 & 0.0355$
  \pm$0.0014 & 3.0 \\
      & & FLTrust & 0.0463$\pm$0.0007 & 0.0519$\pm$0.0090 & 0.0454$\pm$0.0006 & 0.0459$
  \pm$0.0007 & 4.0 \\
      & & FLDetector & 0.0361$\pm$0.0007 & 0.0552$\pm$0.0112 & 0.0304$\pm$0.0007 & 0.0304$
  \pm$0.0007 & 2.0 \\
      & & \cellcolor{mygray}TriGuardFL & \cellcolor{mygray}\textbf{0.0338$\pm$0.0011} &
      \cellcolor{mygray}\textbf{0.0292$\pm$0.0015} &
      \cellcolor{mygray}\textbf{0.0298$\pm$0.0005} &
      \cellcolor{mygray}\textbf{0.0298$\pm$0.0006} & \cellcolor{mygray}1.0 \\
      & & NoAttack+Mean(Base) & 0.0294$\pm$0.0006 & 0.0294$\pm$0.0006 & 0.0294$\pm$0.0006 &
  0.0294$\pm$0.0006 & -- \\
      \bottomrule
      \end{tabular}
  \end{table*}

To assess the scalability of our method, we conducted experiments in more hostile environments by increasing the network scale. This analysis investigates how defense mechanisms perform under heightened threat levels, a critical consideration for real-world deployment.


Fig. \ref{figaccuracymoreclientscifar100} illustrates the evolution of the test accuracy of \texttt{TriGuardFL} with another four other state-of-the-art defense rules  when the numbers of clients are 40 and 60 and propotion of attackers is $20\%$. The accuracy of five defense rules decreases in general. It is obvious that \texttt{TriGuardFL}  achieves near convergence, while NormClipping and FLDetector crash in some cases.

Table \ref{tablosscifar100clientsadv02} compares the test loss of the defense mechanisms in a baseline configuration and a scaled network configuration with a proportionally larger number of attackers. The results highlight the superior scalability of \texttt{TriGuardFL}. As the threat level intensifies, \texttt{TriGuardFL} is the only defense that maintains its top-ranked performance. In contrast, the effectiveness of other defenses degrades sharply; for instance, the test loss for NormClipping and FLDetector increases dramatically, indicating a performance collapse. This demonstrates that \texttt{TriGuardFL}'s detection and reputation mechanisms are robust to scaling, making it a more reliable defense for larger real-world FL systems.

\section{Conclusion}
\label{sec:conclusion}

We presented CARAT, a class-aware robust aggregation rule that validates client updates through hidden task certificates rather than relying only on geometric centrality or a single trusted reference direction. By scoring updates across multiple hidden, class-balanced semantic slices of the task and combining those scores with rank-space anomaly suppression and capped-simplex weight optimization, CARAT is designed to preserve benign heterogeneity while reducing the influence of semantically harmful or structurally extreme clients.

The current repository-backed evidence supports this mechanism claim at a preliminary empirical level: in the completed CIFAR100/FangAttack pilot, CARAT suppresses attacker weights and substantially improves over Mean. At the same time, the present manuscript does not yet support a full benchmark claim, because the clean-accuracy, ASR, non-i.i.d., and multi-seed CARAT suites are still being completed, and the method retains the explicit assumption of server-side labeled reference data. Even with these limitations, the hidden-task certificate view provides a concrete and technically grounded direction for robust FL aggregation, especially in settings where class-selective failures matter.

\bibliographystyle{unsrtnat}
\bibliography{ref}

\newpage
\appendix

\section{Algorithmic Details}
\label{sec:appendix-algorithm}

\subsection{Round-Level Pseudocode}

Algorithm~\ref{alg:carat-round} gives a standard round-level pseudocode view of one CARAT server round.

\begin{algorithm}[t]
\caption{CARAT server aggregation at round $r$}
\label{alg:carat-round}
\begin{algorithmic}[1]
\REQUIRE Current global model $w_r$, client messages $\{u_i^r\}_{i=1}^n$, previous-round weights $\alpha^{\mathrm{prev}}$, probe pool $\mathcal{P}$, CARAT hyperparameters
\ENSURE Aggregated update $\Delta_{\mathrm{agg}}$, updated temporal prior $\alpha^{\mathrm{prev}}$
\IF{$\mathcal{P}$ is uninitialized}
    \STATE Build a balanced probe pool from the server-visible training dataset
\ENDIF
\IF{$r=1$}
    \STATE $\alpha^{\mathrm{prev}} \leftarrow \frac{1}{n}\mathbf{1}$
\ENDIF
\FOR{$i=1$ to $n$}
    \STATE Convert $u_i^r$ into a vector-form delta $\Delta_i^r$
\ENDFOR
\STATE Compute the clipping threshold $\tau_r$ from $\{\|\Delta_i^r\|_2\}_{i=1}^n$
\FOR{$i=1$ to $n$}
    \STATE $\bar{\Delta}_i^r \leftarrow \min\!\left(1,\frac{\tau_r}{\|\Delta_i^r\|_2+\varepsilon}\right)\Delta_i^r$
\ENDFOR
\STATE Choose the evaluation radius $q_r$ from a quantile of positive clipped norms
\FOR{$i=1$ to $n$}
    \STATE $\tilde{\Delta}_i^r \leftarrow q_r \bar{\Delta}_i^r / (\|\bar{\Delta}_i^r\|_2+\varepsilon)$
\ENDFOR
\STATE Sample $T$ class-balanced hidden tasks from $\mathcal{P}$ with replacement/subset fallback if strict balanced sampling fails
\FOR{$t=1$ to $T$}
    \STATE Evaluate the baseline loss $\ell_t(w_r)$
\ENDFOR
\FOR{$i=1$ to $n$}
    \FOR{$t=1$ to $T$}
        \STATE $c_{i,t} \leftarrow \ell_t(w_r) - \ell_t(w_r + \tilde{\Delta}_i^r)$
    \ENDFOR
\ENDFOR
\STATE Robustly standardize certificate columns to obtain $C^{\mathrm{std}}$
\STATE Compute rank penalties $r$ from repeated coordinate subsamples of $\{\bar{\Delta}_i^r\}_{i=1}^n$
\STATE $\alpha^{(0)} \leftarrow \alpha^{\mathrm{prev}}$
\FOR{$k=0$ to $K-1$}
    \STATE $z \leftarrow (C^{\mathrm{std}})^\top \alpha^{(k)}$
    \STATE $\pi \leftarrow \softmin$ task weights induced by $z$ and $\beta$
    \STATE $g \leftarrow C^{\mathrm{std}}\pi - \lambda_r r - 2\lambda_p(\alpha^{(k)}-\alpha^{\mathrm{prev}})$
    \STATE $\alpha^{(k+1)} \leftarrow \Pi_{\mathcal{A}(\hat{\rho})}\!\left(\alpha^{(k)}+\eta g\right)$
\ENDFOR
\STATE $\alpha^\star \leftarrow \alpha^{(K)}$ and $\alpha^{\mathrm{prev}} \leftarrow \alpha^\star$
\STATE $\Delta_{\mathrm{agg}} \leftarrow \sum_{i=1}^n \alpha_i^\star \bar{\Delta}_i^r$
\STATE \textbf{return} $\Delta_{\mathrm{agg}}, \alpha^{\mathrm{prev}}$
\end{algorithmic}
\end{algorithm}

\subsection{Implementation Notes}

\paragraph{Inputs.}
Current global model $w_r$, client messages $\{u_i^r\}_{i=1}^n$, previous-round weights $\alpha^{\mathrm{prev}}$, server-visible training dataset, and CARAT hyperparameters.

\paragraph{Persistent state.}
CARAT maintains two pieces of state across rounds: a balanced probe pool sampled at initialization from the server-visible training dataset, and the previous-round weight vector $\alpha^{\mathrm{prev}}$, which serves as the temporal prior in the current round.

\paragraph{Initialization.}
At the first round, the implementation initializes $\alpha^{\mathrm{prev}}$ to the uniform vector. This makes the initial round unbiased across clients while still allowing later rounds to exploit temporal consistency once CARAT has observed several optimization trajectories.

\paragraph{Reference-pool construction and hidden-task sampling.}
Let $\mathcal{P}_c \subseteq \mathcal{P}$ denote the subset of probe examples with class label $c$. A hidden task is constructed by sampling a small number of examples from each available class pool and concatenating them into one balanced micro-task. If the desired number of samples per class is not available without replacement, CARAT falls back to replacement and, if needed, to a random subset of the probe pool rather than aborting the round. This fail-soft design preserves the intended semantic diversity of the task sampler while ensuring that the server always obtains a well-defined certificate matrix.

\paragraph{Robust score calibration.}
For notation economy, the main text reuses $C$ after standardization. In implementation terms, the raw certificate matrix is first transformed column-wise into
\[
C^{\mathrm{std}}_{i,t}
=
\frac{c_{i,t}-\operatorname{median}_j(c_{j,t})}
{1.4826 \cdot \operatorname{MAD}_j(c_{j,t})+\varepsilon}.
\]
This removes task-specific scale effects while preserving within-task ordering by Proposition~\ref{prop:task-standardization}. Rank penalties are calibrated similarly. For each coordinate subsample $S_s$, let
\[
\rho_i^{(s)}=
\frac{1}{|S_s|}\sum_{j \in S_s}
\left|
\frac{\operatorname{rank}_{i,j}}{n-1} - \frac{1}{2}
\right|.
\]
CARAT then computes a standardized rank score
\[
z_i^{(s)}=
\frac{\rho_i^{(s)}-\operatorname{median}_\ell(\rho_\ell^{(s)})}
{1.4826 \cdot \operatorname{MAD}_\ell(\rho_\ell^{(s)})+\varepsilon},
\]
and aggregates subsamples through
\[
r_i=\max\left(0,\operatorname{median}_s z_i^{(s)}\right).
\]
As a result, a client is penalized only when it is consistently extreme across repeated coordinate views, not because of a single noisy subsample.

\paragraph{Projected optimization details.}
Again writing $C$ for the standardized certificate matrix, CARAT initializes the optimizer with $\alpha^{(0)}=\alpha^{\mathrm{prev}}$ after the first round, and with the uniform vector at initialization. If
\[
z(\alpha)=C^\top \alpha,
\qquad
\pi_t(\alpha)=\frac{\exp(-\beta z_t(\alpha))}{\sum_{\ell=1}^T \exp(-\beta z_\ell(\alpha))},
\]
then the ascent direction used in projected updates is
\[
g(\alpha)=C\pi(\alpha)-\lambda_r r-2\lambda_p(\alpha-\alpha^{\mathrm{prev}}).
\]
With optimizer step size $\eta$, one iteration takes the form
\[
\alpha^{(k+1)}=
\Pi_{\mathcal{A}(\hat{\rho})}
\left(\alpha^{(k)}+\eta g(\alpha^{(k)})\right).
\]
The projection maintains feasibility throughout optimization, so every intermediate iterate already obeys the nonnegativity, unit-sum, and per-client cap constraints described in the main text.

\paragraph{Round invariants and failure handling.}
Two implementation choices are important for robustness. First, the optimizer never leaves the capped simplex, so every intermediate and final weight vector remains feasible. Second, probe-task construction is fail-soft rather than fail-stop: if strict class-balanced sampling is impossible, the code falls back to replacement or to a random subset, ensuring that the round can still proceed with a well-defined certificate matrix.

\paragraph{Outputs.}
Aggregated server update, updated temporal prior $\alpha^{\mathrm{prev}}$, and diagnostic logs including learned weights, task scores, softmin focus, rank penalties, evaluation radius, clipping threshold, and the identities of the hidden probe tasks used in the current round.

\paragraph{Practical diagnostics.}
The most useful round-level diagnostics for debugging and empirical interpretation are the learned weights, the weakest-task values emphasized by the softmin, the clip threshold $\tau_r$, the evaluation radius $q_r$, and the final rank penalties $r_i$. In successful runs, these quantities help distinguish three qualitatively different cases: benign but diverse clients that retain moderate weight, clearly harmful clients that lose certificate support across multiple tasks, and structurally extreme clients whose influence is reduced primarily by the rank penalty.

\section{Theoretical Analysis}
\label{sec:appendix-theory}

This appendix records the main structural properties used to interpret CARAT. These results are not presented as a complete robustness guarantee against arbitrary adaptive attackers. Instead, they clarify why the method decomposes naturally into a normalized evaluation geometry, a task-aware utility objective, and a constrained robust aggregation step.

\begin{proposition}[Feasibility and support of the capped simplex]
\label{prop:capped-simplex}
Let
\[
\mathcal{A}(\hat{\rho})=
\left\{
\alpha \in \R^n:
\alpha_i \ge 0,\;
\sum_{i=1}^n \alpha_i = 1,\;
\alpha_i \le \frac{1}{(1-\hat{\rho})n}
\right\},
\]
with $\hat{\rho}<1$. Then $\mathcal{A}(\hat{\rho})$ is non-empty, and every $\alpha \in \mathcal{A}(\hat{\rho})$ satisfies
\[
|\operatorname{supp}(\alpha)| \ge (1-\hat{\rho})n.
\]
\end{proposition}

\begin{proof}[Proof sketch]
Let $u = 1/((1-\hat{\rho})n)$. Since $\hat{\rho}<1$, we have $u \ge 1/n$, so the uniform weight vector is feasible and the set is non-empty. For the support bound, if only $k$ coordinates of $\alpha$ are positive, then
\[
1=\sum_{i=1}^n \alpha_i \le k u = \frac{k}{(1-\hat{\rho})n},
\]
which implies $k \ge (1-\hat{\rho})n$. Therefore CARAT can suppress suspicious clients, but cannot collapse the aggregate onto an arbitrarily small subset of participants.
\end{proof}

\begin{proposition}[Softmin bounds and gradient form]
\label{prop:softmin}
Let $s \in \R^T$ and define
\[
\softmin_{\beta}(s)
= -\frac{1}{\beta}\log\left(\frac{1}{T}\sum_{t=1}^T e^{-\beta s_t}\right).
\]
Then
\[
\min_t s_t \le \softmin_{\beta}(s) \le \min_t s_t + \frac{\log T}{\beta}.
\]
Moreover,
\[
\lim_{\beta \to \infty}\softmin_{\beta}(s)=\min_t s_t,
\qquad
\lim_{\beta \to 0}\softmin_{\beta}(s)=\frac{1}{T}\sum_{t=1}^T s_t.
\]
Moreover, if $z(\alpha)=C^\top\alpha$ and
\[
\pi_t(\alpha)=\frac{\exp(-\beta z_t(\alpha))}{\sum_{\ell=1}^T \exp(-\beta z_\ell(\alpha))},
\]
then
\[
\nabla_\alpha \softmin_{\beta}(C^\top\alpha)=C\pi(\alpha).
\]
\end{proposition}

\begin{proof}[Proof sketch]
Let $m=\min_t s_t$. Since
\[
e^{-\beta m}/T \le \frac{1}{T}\sum_{t=1}^T e^{-\beta s_t} \le e^{-\beta m},
\]
applying $-(1/\beta)\log(\cdot)$ yields
\[
m \le \softmin_{\beta}(s) \le m+\frac{\log T}{\beta}.
\]
The limit $\beta \to \infty$ follows immediately from the vanishing gap bound. For $\beta \to 0$, a first-order expansion of the exponential average gives
\[
\frac{1}{T}\sum_{t=1}^T e^{-\beta s_t}
=
1-\beta \frac{1}{T}\sum_{t=1}^T s_t + o(\beta),
\]
so the softmin converges to the arithmetic mean. Differentiating the softmin expression with respect to $z_t$ yields the Gibbs weights $\pi_t(\alpha)$, and the chain rule then gives $C\pi(\alpha)$. Hence the ascent direction is a convex combination of certificate columns, with more weight placed on tasks whose current aggregate utility is smaller.
\end{proof}

\begin{proposition}[Directional invariance of certificate scoring]
\label{prop:directional-invariance}
Suppose two nonzero clipped updates satisfy $\bar{\Delta}_i = a \bar{\Delta}_j$ for some scalar $a>0$. Then their normalized versions in Eq.~\eqref{eq:normalize} coincide, i.e.,
\[
\tilde{\Delta}_i=\tilde{\Delta}_j.
\]
Consequently, their hidden-task certificates are identical:
\[
c_{i,t}=c_{j,t}\qquad \text{for all } t \in \{1,\dots,T\}.
\]
\end{proposition}

\begin{proof}[Proof sketch]
Eq.~\eqref{eq:normalize} rescales every nonzero clipped update to the common radius $q_r$. Positive collinearity is therefore removed by normalization:
\[
q_r\frac{\bar{\Delta}_i}{\|\bar{\Delta}_i\|_2}
=
q_r\frac{a\bar{\Delta}_j}{\|a\bar{\Delta}_j\|_2}
=
q_r\frac{\bar{\Delta}_j}{\|\bar{\Delta}_j\|_2}.
\]
Substituting the resulting equality into Eq.~\eqref{eq:cert} gives identical task scores. Thus, after clipping, the certificate stage depends on direction rather than residual norm magnitude.
\end{proof}

\begin{proposition}[Task-wise standardization preserves within-task ordering]
\label{prop:task-standardization}
Fix a task $t$ and define a standardized certificate score
\[
\hat{c}_{i,t}=\frac{c_{i,t}-m_t}{s_t+\varepsilon},
\]
where $m_t \in \R$ and $s_t>0$ depend only on task $t$ and not on client $i$. Then for any clients $i,j$,
\[
c_{i,t} \ge c_{j,t}
\quad \Longleftrightarrow \quad
\hat{c}_{i,t} \ge \hat{c}_{j,t}.
\]
Moreover,
\[
\hat{c}_{i,t}-\hat{c}_{j,t}
=
\frac{c_{i,t}-c_{j,t}}{s_t+\varepsilon}.
\]
\end{proposition}

\begin{proof}[Proof sketch]
For a fixed task $t$, the map $x \mapsto (x-m_t)/(s_t+\varepsilon)$ is affine with strictly positive slope because $s_t+\varepsilon>0$. Hence it preserves ordering. The difference identity follows by direct subtraction. Therefore column-wise robust standardization calibrates cross-task scale without changing the within-task ranking induced by the raw certificate scores.
\end{proof}

\begin{proposition}[Temporal prior and projected-iterate invariants]
\label{prop:temporal-prior}
Define
\[
f(\alpha)=\softmin_{\beta}(C^\top\alpha)-\lambda_r r^\top\alpha-\lambda_p\|\alpha-\alpha^{\mathrm{prev}}\|_2^2.
\]
If $\alpha,\alpha' \in \mathcal{A}(\hat{\rho})$ satisfy $C^\top\alpha=C^\top\alpha'$ and $r^\top\alpha=r^\top\alpha'$, then
\[
f(\alpha)-f(\alpha')
=
-\lambda_p\!\left(\|\alpha-\alpha^{\mathrm{prev}}\|_2^2-\|\alpha'-\alpha^{\mathrm{prev}}\|_2^2\right).
\]
In addition, under projected gradient ascent
\[
\alpha^{(k+1)}=\Pi_{\mathcal{A}(\hat{\rho})}\!\left(\alpha^{(k)}+\eta \nabla f(\alpha^{(k)})\right),
\]
every iterate remains in $\mathcal{A}(\hat{\rho})$.
\end{proposition}

\begin{proof}[Proof sketch]
The first statement follows by direct substitution: if the certificate and rank terms match, the only remaining difference in objective value is the quadratic prior term. CARAT therefore prefers, among utility-equivalent feasible weights, the solution closer to $\alpha^{\mathrm{prev}}$. The second statement is immediate from the projection operator. Hence the optimizer preserves three round invariants throughout: $\alpha_i \ge 0$, $\sum_i \alpha_i =1$, and $\alpha_i \le 1/((1-\hat{\rho})n)$.
\end{proof}

\begin{proposition}[Norm bound for the final aggregate]
\label{prop:aggregate-bound}
Suppose each clipped update satisfies $\|\bar{\Delta}_i\|_2 \le \tau_r$ and $\alpha \in \mathcal{A}(\hat{\rho})$. Then the final aggregate
\[
\Delta_{\mathrm{agg}}=\sum_{i=1}^n \alpha_i \bar{\Delta}_i
\]
satisfies
\[
\|\Delta_{\mathrm{agg}}\|_2 \le \tau_r.
\]
\end{proposition}

\begin{proof}[Proof sketch]
Because $\alpha_i \ge 0$ and $\sum_i \alpha_i=1$, the aggregate is a convex combination of clipped updates. By the triangle inequality,
\[
\|\Delta_{\mathrm{agg}}\|_2
\le
\sum_{i=1}^n \alpha_i \|\bar{\Delta}_i\|_2
\le
\sum_{i=1}^n \alpha_i \tau_r
=
\tau_r.
\]
Thus the final server step inherits the clipping radius as a hard norm upper bound.
\end{proof}

\begin{proposition}[Per-round complexity]
\label{prop:complexity}
Let $n$ be the number of clients, $T$ the number of hidden tasks, $m$ the number of sampled coordinates for rank scoring, $S$ the number of rank subsamples, and $K$ the number of projected-gradient steps. Ignoring constant factors from batch size and model width, the extra defense-side work of CARAT scales as
\[
\text{probe-model forward passes} + O(Snm + KnT).
\]
\end{proposition}

\begin{proof}[Proof sketch]
CARAT performs one baseline probe evaluation plus $n$ candidate probe evaluations over the $T$ hidden tasks, then computes $S$ rank summaries over $m$ coordinates, and finally executes $K$ projected optimization steps over $n$ weights and $T$ task scores. In the current implementation, the probe-model forward passes dominate wall-clock time, while the rank and optimization stages are lower-order terms controlled by $m$, $S$, and $K$.
\end{proof}

\paragraph{Remark on assumptions.}
The theoretical object $D^{\mathrm{ref}}$ in the main text should be read as the probe pool \emph{after} subsampling. The code, however, obtains that pool by sampling from the full server-visible training dataset. Any final formal discussion should keep this distinction explicit, because it materially strengthens the server-side assumption relative to a fully data-free defense.

\section{Additional Experimental Details}
\label{sec:appendix-experiments}

This section summarizes the current verified experimental scope together with the next empirical extensions already supported by the repository structure.

\paragraph{Currently verified run.}
The strongest directly verified completed CARAT run in the repository uses FedAvg, CIFAR100, ResNet-18, an i.i.d.\ partition, 20 clients, 4 attackers, FangAttack, 300 rounds, learning rate $0.05$, local epochs $4$, and \verb|evaluate=false|. The main text reports the final training metrics and runtime of this run together with the corresponding Mean and TriGuardFL baselines. In the current repository snapshot dated April 16, 2026, the evaluation-enabled \texttt{main\_untargeted.yaml} sweep has also started producing ALIE baseline logs, but no CARAT task from that sweep is complete yet.

\paragraph{Supported benchmark matrix.}
The committed \texttt{main\_untargeted.yaml} spec already expands to 216 tasks over i.i.d.\ and non-i.i.d.\ splits with Dirichlet $\alpha \in \{0.1,0.5\}$, attacker fractions $\{0.2,0.4\}$, untargeted attacks \{ALIE, MinMax, FangAttack\}, defenses \{Mean, FLTrust, MultiKrum, CARAT\}, and three repeats. Companion specs \texttt{clean\_reference.yaml}, \texttt{backdoor.yaml}, and \texttt{ablations.yaml} cover clean baselines, backdoor evaluation, and component analysis.

\paragraph{Next additions.}
\begin{enumerate}
\item Complete the running \texttt{main\_untargeted.yaml} sweep, prioritizing the CARAT rows and reporting multi-seed means and variances for clean test accuracy and loss.
\item Extend the verified runs from the current i.i.d.\ slice to the non-i.i.d.\ Dirichlet settings already present in \texttt{main\_untargeted.yaml}.
\item Complete the committed \texttt{backdoor.yaml} suite so that the paper reports both clean utility and ASR.
\item Run the config-supported ablations in Table~\ref{tab:ablation-placeholder}, prioritizing the number of hidden tasks, the rank and prior weights, and the clipping controls.
\item Report multi-seed means and variances, plus defense runtime overhead, before making any strong comparative claim against other robust aggregators.
\end{enumerate}

\end{document}
